% An appendix
%======================================================================
\chapter{Wybrane Fragmenty Kodu Źródłowego}
\label{ch:Appendix-Code}
%======================================================================

Listingi w tym dodatku to fragmenty mojego kodu wybrane pod kątem obrony. Każdy z nich pokazuje konkretny, samodzielny kawałek logiki, do którego odsyła główna część pracy - od maszyny stanów, przez śmierć wroga, po serce \texttt{EliteBrain}. Pełne pliki znajdują się w repozytorium projektu, ale aby zachować czytelność na papierze, w niektórych miejscach skróciłem mniej istotne gałęzie kodu (komentarz \texttt{// ...} oznacza takie skrócenie).

\section{Maszyna Stanów}

Poniższy kod przedstawia implementację kontrolera maszyny stanów oraz klasę bazową stanu postaci.

\begin{lstlisting}[caption={StateMachine.cs oraz CharacterState.cs},label={lst:sm}]
// StateMachine.cs
public partial class StateMachine : Node
{
    private CharacterState _currentState;

    public void SwitchState<T>() where T : CharacterState
    {
        _currentState?.ExitState();
        _currentState = GetOrCreate<T>();
        _currentState.EnterState();
    }

    public override void _Process(double delta)
        => _currentState?.UpdateState(delta);

    public override void _PhysicsProcess(double delta)
        => _currentState?.PhysicsUpdateState(delta);
}

// CharacterState.cs - klasa bazowa stanow
public abstract partial class CharacterState : Node
{
    protected Character characterNode;

    public override void _Ready()
        => characterNode = GetOwner<Character>();

    protected virtual void EnterState()  { }
    protected virtual void ExitState()   { }
    public virtual void UpdateState(double delta)       { }
    public virtual void PhysicsUpdateState(double delta){ }
}
\end{lstlisting}

\section{Stan Śmierci Przeciwnika}

Fragment implementujący wieloetapową logikę śmierci wroga z obsługą questów i animacji.

\begin{lstlisting}[caption={EnemyDeathState.cs - wieloetapowa smierc wroga},label={lst:death}]
public partial class EnemyDeathState : EnemyState
{
    private bool _isDying = false;

    protected override void EnterState()
    {
        if (_isDying) return;
        _isDying = true;

        if (characterNode is Enemy enemy)
        {
            enemy.SpawnDrops();
            if (!string.IsNullOrEmpty(enemy.EnemyId))
                GameEvents.RaiseEnemyKilled(enemy.EnemyId);
        }

        characterNode.SetCollisionLayerValue(3, false);
        characterNode.SetCollisionMaskValue(2, false);

        if (characterNode.AnimationPlayerNode.HasAnimation("Death"))
        {
            characterNode.AnimationPlayerNode
                .AnimationFinished += HandleAnimationFinished;
            characterNode.AnimationPlayerNode.Play("Death");
        }
        else
        {
            RemoveEnemy();
        }
    }

    private void RemoveEnemy()
    {
        if (characterNode.PathNode != null)
            characterNode.PathNode.QueueFree();
        else
            characterNode.QueueFree();
    }
}
\end{lstlisting}

\section{QuestManager - Rejestrowanie Zabójstw}

Fragment systemu questów odpowiedzialny za śledzenie postępu celu Kill. Pełna wersja obsługuje także questy sekwencyjne (pomija cele poza aktywnym indeksem).

\begin{lstlisting}[caption={QuestManager.cs - RegisterKill},label={lst:qm-kill}]
public void RegisterKill(string enemyId)
{
    var questsToComplete = new List<QuestResource>();

    foreach (var quest in ActiveQuests)
    {
        if (quest.Objectives == null) continue;

        for (int i = 0; i < quest.Objectives.Length; i++)
        {
            var obj = quest.Objectives[i];
            if (obj == null) continue;
            if (obj.Type != QuestObjectiveType.Kill) continue;
            if (obj.TargetId != enemyId) continue;

            // Questy sekwencyjne: aktualizuj tylko aktywny cel
            if (quest.SequentialObjectives
                && _seqNextIndex.TryGetValue(quest.Id, out int idx)
                && idx != i) continue;

            if (!_killProgress.ContainsKey(quest.Id))
                _killProgress[quest.Id] = new Dictionary<string, int>();
            if (!_killProgress[quest.Id].ContainsKey(enemyId))
                _killProgress[quest.Id][enemyId] = 0;

            _killProgress[quest.Id][enemyId]++;

            if (_killProgress[quest.Id][enemyId] >= obj.RequiredAmount)
            {
                if (quest.SequentialObjectives)
                    AdvanceSequentialObjective(quest);
                else if (!questsToComplete.Contains(quest))
                    questsToComplete.Add(quest);
            }
        }
    }

    foreach (var quest in questsToComplete)
        CompleteQuest(quest);
}
\end{lstlisting}

\section{EliteBrain - Bandyta Kontekstowy}

\subsection{Wybór Taktyki ($\varepsilon$-greedy)}

\begin{lstlisting}[caption={EliteBrain.SelectTactic},label={lst:eb-select}]
// Wybor taktyki: e-greedy z uwzglednieniem kontekstu (HP, liczba elit)
private Tactic SelectTactic(int hpBucket, int eliteBucket)
{
    if (_rng.Randf() < Epsilon)
        return (Tactic)_rng.RandiRange(0, TacticCount - 1);

    int best = 0;
    float bestVal = QTable[0, hpBucket, eliteBucket];
    for (int i = 1; i < TacticCount; i++)
    {
        float v = QTable[i, hpBucket, eliteBucket];
        if (v > bestVal)
        {
            bestVal = v;
            best = i;
        }
    }
    return (Tactic)best;
}
\end{lstlisting}

\subsection{Aktualizacja Wartości Akcji (Q-learning EMA)}

\begin{lstlisting}[caption={EliteBrain.UpdateTacticValue},label={lst:eb-update}]
private void UpdateTacticValue(Tactic t, int hpBucket, int eliteBucket, float reward)
{
    int idx = (int)t;
    NTable[idx, hpBucket, eliteBucket]++;
    float oldQ = QTable[idx, hpBucket, eliteBucket];
    QTable[idx, hpBucket, eliteBucket] =
        oldQ + LearningRate * (reward - oldQ);
}
\end{lstlisting}

\subsection{Dyskretyzacja Kontekstu}

\begin{lstlisting}[caption={EliteBrain.GetCurrentContext - kubelki HP i liczby elit},label={lst:eb-ctx}]
public (int hp, int elites) GetCurrentContext()
{
    int hp = 2; // High
    var player = Player.Instance;
    if (player != null)
    {
        var hpStat = player.GetStatResource(Stat.Health);
        if (hpStat != null)
        {
            if (hpStat.StatValue > MaxObservedPlayerHp)
                MaxObservedPlayerHp = hpStat.StatValue;

            float ratio = MaxObservedPlayerHp > 0.001f
                ? hpStat.StatValue / MaxObservedPlayerHp : 1f;

            if (ratio < 0.33f)      hp = 0; // Low
            else if (ratio < 0.66f) hp = 1; // Mid
            else                    hp = 2; // High
        }
    }

    int elites;
    int alive = GetTree().GetNodesInGroup("elite_enemies").Count;
    if (alive <= 1)      elites = 0; // 1 elita
    else if (alive == 2) elites = 1; // 2 elity
    else                 elites = 2; // 3+ elit

    return (hp, elites);
}
\end{lstlisting}

\subsection{Cykl Życia Starcia}

\begin{lstlisting}[caption={EliteBrain.StartEncounter / EndEncounter},label={lst:eb-encounter}]
private void StartEncounter()
{
    _encounterActive   = true;
    _encounterReward   = 0f;
    _encounterDuration = 0f;

    (CurrentCtxHp, CurrentCtxElites) = GetCurrentContext();
    CurrentTactic = SelectTactic(CurrentCtxHp, CurrentCtxElites);
}

private void EndEncounter()
{
    _encounterActive = false;

    float norm = Mathf.Clamp(
        _encounterReward / Mathf.Max(1.5f, _encounterDuration),
        -5f, 5f);

    UpdateTacticValue(CurrentTactic,
                      CurrentCtxHp, CurrentCtxElites, norm);

    Epsilon = Mathf.Max(EpsilonFloor, Epsilon * EpsilonDecay);
    AdaptParametersFromPlayerModel();
    SaveBrain();
}
\end{lstlisting}

\subsection{Predykcja Pozycji Gracza i Cel Ruchu Elity}

\begin{lstlisting}[caption={EliteBrain.ComputeChaseDestination (fragment)},label={lst:eb-chase}]
public Vector3 ComputeChaseDestination(int index, int total, Vector3 elitePos)
{
    var player = Player.Instance;
    if (player == null) return elitePos;

    Vector3 pPos = player.GlobalPosition;
    Vector3 pVel = _havePrevPlayer
        ? (pPos - _prevPlayerPos) / Mathf.Max(0.016f, (float)_lastDelta)
        : Vector3.Zero;

    // lead-shot: przewidywana pozycja za LeadTime sekund
    Vector3 predicted = pPos + LeadTime * pVel;

    switch (CurrentTactic)
    {
        case Tactic.Rush:
            return predicted;

        case Tactic.Kite:
        {
            Vector3 away = (elitePos - predicted).Normalized();
            return predicted + PreferredDistance * away;
        }

        case Tactic.Flank:
        {
            float theta = (index * Mathf.Tau) / Mathf.Max(1, total);
            Vector3 offset = new Vector3(
                Mathf.Cos(theta), 0, Mathf.Sin(theta))
                * PreferredDistance * 0.7f;
            return predicted + offset;
        }

        default: // Ambush - wolniejsza aktualizacja celu
            if ((predicted - _lastAmbushTarget).Length() > 3f)
                _lastAmbushTarget = predicted;
            return _lastAmbushTarget;
    }
}
\end{lstlisting}

\section{Zapis Stanu Mózgu AI}

\begin{lstlisting}[caption={EliteBrain.SaveBrain (fragment) - serializacja do JSON},label={lst:eb-save}]
public void SaveBrain()
{
    var dto = new EliteBrainDto
    {
        SchemaVersion       = 2,
        TacticCount         = TacticCount,
        HpBuckets           = HpBuckets,
        EliteBuckets        = EliteBuckets,
        QFlat               = FlattenFloats(QTable),
        NFlat               = FlattenInts(NTable),
        Epsilon             = Epsilon,
        Aggression          = Aggression,
        Caution             = Caution,
        PreferredDistance   = PreferredDistance,
        LeadTime            = LeadTime,
        MaxObservedPlayerHp = MaxObservedPlayerHp,
    };

    string json = JsonSerializer.Serialize(dto,
        new JsonSerializerOptions { WriteIndented = true });

    using var f = FileAccess.Open(BrainPath, FileAccess.ModeFlags.Write);
    f?.StoreString(json);
}
\end{lstlisting}

\section{System Dropów Przeciwnika}

\begin{lstlisting}[caption={Enemy.SpawnDrops - losowanie z Fisher-Yates},label={lst:drops}]
public void SpawnDrops()
{
    if (DropOrbScene == null || Drops == null || Drops.Length == 0)
        return;

    var rng = new RandomNumberGenerator();
    rng.Randomize();

    var potions = new Array<ItemResource>();
    var others  = new Array<ItemResource>();

    foreach (var it in Drops)
    {
        if (it == null) continue;
        if (it.IsConsumable) potions.Add(it);
        else                 others.Add(it);
    }

    var droppedOther  = RollOneFromList(others,  rng);
    var droppedPotion = RollOneFromList(potions, rng);

    if (droppedOther  != null) SpawnOneOrb(sceneRoot, droppedOther,  rng);
    if (droppedPotion != null) SpawnOneOrb(sceneRoot, droppedPotion, rng);
}

// Tasowanie Fisher-Yates + sprawdzenie DropChance
private ItemResource RollOneFromList(Array<ItemResource> list,
                                     RandomNumberGenerator rng)
{
    for (int i = list.Count - 1; i > 0; i--)
    {
        int j = rng.RandiRange(0, i);
        (list[i], list[j]) = (list[j], list[i]);
    }
    foreach (var it in list)
        if (rng.Randf() <= it.DropChance)
            return it;
    return null;
}
\end{lstlisting}

\section{Globalny System Zdarzeń}

\begin{lstlisting}[caption={GameEvents.cs - wybrane zdarzenia},label={lst:events}]
public class GameEvents
{
    public static event Action<ItemResource, int> OnItemCollected;
    public static event Action<RewardResource>    OnReward;
    public static event Action<int>               OnNewEnemyCount;
    public static event Action<string>            OnEnemyKilled;
    public static event Action<float>             OnDamageDealt;
    public static event Action<float>             OnDamageTaken;
    public static event Action                    OnDungeonEntered;
    public static event Action                    OnDungeonExited;
    public static event Action<string>            OnJournalEntry;
    public static event Action<string, float>     OnToast;
    public static event Action<float>             OnCameraShake;

    public static void RaiseToast(string msg, float seconds = 2.5f)
        => OnToast?.Invoke(msg, seconds);
    public static void RaiseJournalEntry(string text)
        => OnJournalEntry?.Invoke(text);
    public static void RaiseEnemyKilled(string id)
        => OnEnemyKilled?.Invoke(id);
    // ... pozostale Raise-ery analogicznie
}
\end{lstlisting}

\section{Shader Postaci - Tint, Flash i Niewidoczność}

Pełny shader \texttt{character\_\allowbreak hurt.\allowbreak gdshader} obsługuje trzy zachowania: flash trafienia, tint elit i bossa oraz niewidoczność niewidzialnych przeciwników. Ostatni uniform (\texttt{sprite\_\allowbreak visibility}) jest tu kluczowy - \texttt{Sprite3D.\allowbreak Modulate} nie wystarcza w trybie \texttt{unshaded}, bo overlay czyta teksturę niezależnie od \texttt{Modulate}. Instrukcja \texttt{discard} eliminuje fragment z obrazu wynikowego, a nie tylko go przezroczyśli.

\begin{lstlisting}[caption={Shaders/character\_hurt.gdshader - fragment() z tintem i niewidocznoscia},label={lst:shader},language={}]
shader_type spatial;
render_mode unshaded, cull_disabled, depth_draw_opaque;

uniform sampler2D tex;
uniform float     hurt_flash         : hint_range(0, 1) = 0.0;
uniform float     elite_tint_strength = 0.0;
uniform vec3      elite_tint_color    = vec3(0.82, 0.58, 1.12);
// 0 = niewidzialny sprite (np. InvisibleEnemy), 1 = pelna widocznosc
uniform float     sprite_visibility   = 1.0;

void fragment() {
    vec4 color = texture(tex, UV);
    ALPHA = color.a * sprite_visibility;

    // sprite_visibility blisko zera -> calkowicie usun fragment
    if (sprite_visibility < 0.01) {
        discard;
    }

    // Tint elit / bossa (ten sam mechanizm, inna paleta)
    if (elite_tint_strength > 0.001) {
        vec3 tinted = color.rgb * elite_tint_color;
        color.rgb  = mix(color.rgb, tinted, elite_tint_strength);
    }

    // Flash czerwieni po trafieniu
    color.rgb = mix(color.rgb, vec3(1.0, 0.15, 0.15), hurt_flash);

    ALBEDO = color.rgb;
}
\end{lstlisting}

\section{Radar - Detekcja Niewidzialnych w Półkuli}

Fragment radaru, w którym co klatkę fizyki zbierane są ciała w \texttt{Area3D}, filtrowane do \texttt{InvisibleEnemy} i sprawdzane warunkiem górnej półkuli. Sygnały \texttt{BodyEntered}/\texttt{BodyExited} przy szybkim przekraczaniu progu okazały się zawodne, stąd przejście na \texttt{Get\allowbreak Overlapping\allowbreak Bodies()} z osobnym \texttt{HashSet}-em ,,kto był w środku w poprzedniej klatce".

\begin{lstlisting}[caption={Radar.cs - akumulacja czasu w polkuli (skrocone)},label={lst:radar}]
public override void _PhysicsProcess(double delta)
{
    if (RotatingDish != null)
        RotatingDish.RotateY((float)delta * DishRotationSpeed);

    if (DetectionArea == null) return;

    var insideNow = new HashSet<InvisibleEnemy>();

    foreach (var body in DetectionArea.GetOverlappingBodies())
    {
        if (body is not InvisibleEnemy inv || inv.IsRevealed) continue;
        if (!IsInsideHemisphere(inv.GlobalPosition)) continue;

        insideNow.Add(inv);
        inv.AccumulateRadarTime(delta);  // wewn. trigger Reveal() po 3 s
    }

    // ci, ktorzy w tej klatce wyszli z polkuli - reset
    foreach (var inv in _insideLastFrame)
        if (!insideNow.Contains(inv))
            inv.ResetRadarAccumulation();

    _insideLastFrame.Clear();
    foreach (var inv in insideNow)
        _insideLastFrame.Add(inv);
}

private bool IsInsideHemisphere(Vector3 worldPos)
{
    Vector3 local = worldPos - GlobalPosition;
    if (local.LengthSquared() > DetectionRadius * DetectionRadius) return false;
    return local.Y >= -0.35f;  // gorna polowa + tolerancja na nierownosci podlogi
}
\end{lstlisting}

\section{Minimapa - Maska Kołowa w \texttt{canvas\_item}}

Mapa renderuje się w \texttt{SubViewport} i przechodzi przez prosty shader na \texttt{TextureRect}-ie wyświetlającym tę teksturę. \texttt{smoothstep} daje miękką krawędź, a wartości progowe dobrałem tak, żeby zawartość pokrywała się z wewnętrznym pierścieniem ramki minimapy.

\begin{lstlisting}[caption={Shaders/minimap\_mask.gdshader - maska kolowa},label={lst:mini-mask},language={}]
shader_type canvas_item;

uniform sampler2D source;

void fragment() {
    vec4 c = texture(source, UV);
    float d = distance(UV, vec2(0.5));
    float mask = 1.0 - smoothstep(0.395, 0.412, d);
    COLOR = vec4(c.rgb, c.a * mask);
}
\end{lstlisting}

\section{Serializacja Autosave'a}

\begin{lstlisting}[caption={SaveGameService - struktury DTO dla autosave},label={lst:save}]
public class AutosaveDto
{
    public int   SaveVersion         { get; set; }
    public float Px { get; set; }
    public float Py { get; set; }
    public float Pz { get; set; }
    public float BaseHealth          { get; set; }
    public float BaseStrength        { get; set; }
    public int   Gold                { get; set; }
    public int   Experience          { get; set; }
    public List<InvSlotDto>          Inventory        { get; set; } = new();
    public List<EquipSlotSaveDto>    EquipmentSlots   { get; set; }
    public List<string>              ActiveQuestPaths { get; set; } = new();
    public List<string>              CompletedQuestPaths { get; set; } = new();
    public List<KillProgDto>         KillProgress     { get; set; } = new();
    public List<QuestSeqProgDto>     QuestSeqProgress { get; set; } = new();
    public int  DungeonRunsCompleted   { get; set; }
    public bool DungeonRelicEverGranted{ get; set; }
}

public class InvSlotDto
{
    public string Path              { get; set; }
    public int    Count             { get; set; }
    public int    ItemUpgradeLevel  { get; set; }
    public string ItemDisplayName   { get; set; }
    public int    ItemAttackBonus   { get; set; }
    public int    ItemDefenseBonus  { get; set; }
    public int    ItemSpeedBonus    { get; set; }
}
\end{lstlisting}
